% Generated draft: requires mathematical and source review.
\documentclass[11pt]{article}
\usepackage[margin=25mm]{geometry}
\usepackage{fontspec}
\setmainfont{Latin Modern Roman}
\usepackage{amsmath,amssymb,mathtools,unicode-math}
\setmathfont{Latin Modern Math}
\usepackage{enumitem,needspace,xurl,hyperref}
\setlength{\parindent}{0pt}
\setlength{\parskip}{5pt}
\setlength{\emergencystretch}{3em}
\newcommand{\N}{\mathbb{N}}
\newcommand{\Z}{\mathbb{Z}}
\newcommand{\Q}{\mathbb{Q}}
\newcommand{\R}{\mathbb{R}}
\newcommand{\C}{\mathbb{C}}
\newcommand{\F}{\mathbb{F}}
\newcommand{\A}{\mathbb{A}}
\newcommand{\PP}{\mathbb{P}}
\newcommand{\E}{\mathbb{E}}
\newcommand{\eps}{\varepsilon}
\newcommand{\dd}{\,\mathrm{d}}
\begin{document}
\section*{UnsolvedMath ID: 30002751}
\textbf{Source problem code:} catalogue entry
\subsection{Definitions and mathematical statement}
TASEP denotes the totally asymmetric simple exclusion process: particles obey an exclusion constraint (at most one particle per site), and allowed particle jumps have a single preferred direction. This terminology alone does not specify the driven process intended in the supplied record. In particular, the underlying sites, transition rates, driving mechanism, boundary rules, and the meaning and admissible values of the parameter \(\alpha\) are not supplied.

The source asks: \emph{Can the different behavior of the boundary conditions, as a function of \(\alpha\) in the driven TASEP, be understood through a large-deviation argument?}

The requested mathematical task is to derive a large-deviation explanation of the boundary behavior referred to in that question. However, the record does not describe that behavior, specify the relevant observable or probability laws, or identify a limiting regime and large-deviation speed. Consequently, it does not determine a precise proposition with explicit hypotheses, parameter ranges, or quantifiers. A definite mathematical formulation requires the missing model and boundary-behavior specifications; none is selected here.
\subsection{Short English statement}
The question asks whether large-deviation theory explains how boundary behavior in a driven TASEP depends on \(\alpha\). The supplied record leaves both \(\alpha\) and the boundary behavior unspecified, so a precise conjecture cannot yet be extracted. Its open status is a supplied curation label.
\subsection{Sources}
\begin{itemize}
\item Local UnsolvedMath registry, supplied complete record for permanent ID \(30002751\), source problem code \texttt{catalogue entry}, title \emph{Large-Deviation Explanation of Driven TASEP Boundary Behavior}; collection: Oberwolfach Reports Open Problems. No individual proposer is identified.
\item Supplied original report citation: \emph{Variational Methods for Evolution (2015)}, Oberwolfach Reports, pp.~3177--3254, DOI \texttt{10.4171/owr/2014/57}. \url{https://doi.org/10.4171/owr/2014/57}. The background describes Oberwolfach Reports as documenting mathematical workshops and their open-problem discussions. The record gives proposed year 2014.
\item Supplied literature source labeled ``current open ASEP large-deviation analysis'': \url{https://doi.org/10.1007/s00440-026-01486-w}. No authors, article title, or specific supporting result are supplied.
\item Literature assessment checked at \texttt{2026-08-21}: the supplied curation labels the problem open and states that large-deviation analyses give model-specific rate functions and phase diagrams for boundary-driven systems, but that no general derivation of the reported \(\alpha\)-dependent boundary behavior has been established. This is supplied dated literature triage derived from the OWR curation, not an independently verified assessment. Its time-sensitive status and cited claims require independent verification. Both literature-source links above are preserved as supplied.
\end{itemize}
\end{document}
